% =====================================================================
% Caption for arg_vs_ml_consequences_rwp_v8.png, and the replacement for the
% figure-description paragraph (paragraph 2 of the subsection).
%
% The panel inventory changed with v8: surface precipitation rate and in-cloud
% cloud water are gone; rain water path and cloudy-column fraction are each
% split by background.  The old paragraph 2 and caption describe panels that
% no longer exist, so both must be replaced even though the surrounding prose
% is unchanged.
% =====================================================================

% ---- replacement for paragraph 2 of the subsection --------------------
Figure~\ref{fig:arg_consequences} shows time series of aerosol- and cloud-related variables from $09{:}00$ to $13{:}00$~UTC on 15~July. Aerosol and droplet quantities (Figure~\ref{fig:arg_consequences}a--e) are air-mass-weighted means below $2000~\mathrm{m}$ above ground level (AGL) over cloudy columns ($q_c > 10^{-6}~\mathrm{kg\,kg^{-1}}$). The liquid and rain water paths (Figure~\ref{fig:arg_consequences}f, g, i) are full-column integrals area-averaged over the same cloudy columns, the full column being necessary for rain, which falls below cloud base. Cloudy-column fraction (Figure~\ref{fig:arg_consequences}h, j) is the fraction of the domain meeting that cloud test. Rain water path and cloud fraction are each shown twice, once for the moist pair and once for the default-humidity pair, because the two backgrounds differ by two to three orders of magnitude in rain water and would otherwise share an axis on which the weaker pair is unreadable. The droplet mean-volume radius, $r_v = \left(3 q_c / 4\pi\rho_w N_d\right)^{1/3}$, is formed from the ratio of the averaged cloud water to the averaged droplet number. Cloud condensation nuclei (CCN) are evaluated at $S=0.05\%$.

% ---- figure caption ---------------------------------------------------
\begin{figure}
\centering
\includegraphics[width=\textwidth]{Figures/arg_vs_ml_consequences_rwp_v8.png}
\caption{Consequences of the ARG activation bias for the simulated aerosol and cloud, $09{:}00$--$13{:}00$~UTC on 15~July 2017. Solid curves are \texttt{ML-ACT}, dashed \texttt{ARG-ACT}; colours identify the four cases (blue \texttt{org}, red \texttt{org\_3aer}, green \texttt{mid}, purple \texttt{mid\_3aer}). (a) accumulation-mode and (b) Aitken-mode aerosol number, (c) CCN at $S=0.05\%$, (d) cloud droplet number and (e) droplet mean-volume radius, all air-mass-weighted means below $2000~\mathrm{m}$ AGL over cloudy columns; (f) liquid water path, a full-column integral area-averaged over cloudy columns. Rain water path and cloudy-column fraction are split by background so that each pair keeps a readable scale: (g) rain water path and (h) cloudy-column fraction for the moist cases, (i) and (j) the same two quantities for the default-humidity cases. Rain water path is plotted linearly for the moist pair, which spans a factor of $2.5$, and logarithmically for the default-humidity pair, which spans a factor of $289$. Note that the vertical axis of panel~(h) does not include zero: the moist deck is overcast throughout and its cloud fraction varies only between $0.94$ and $0.98$, so the excursions there are two orders of magnitude smaller than those in panel~(j) despite occupying comparable space on the page. All eight runs branch from a common restart at $09{:}02$~UTC and differ only in which activated fraction is passed to the microphysics.}
\label{fig:arg_consequences}
\end{figure}
